\scanpage{76}\mspage{75}

Dans ce cas, l'axiome $C6$ (pour une relation \inserted{d'ordre} \unsure{suppl|suppl.} $\preceq$ sur $\unsure{\mathcal F|\mathscr F}=\unsure{\mathrm P|\mathcal P}_{f}(\unsure{M|\mathcal M})$) signifie \unsure{seulement|seulmt} que si $F'\preceq F$, $F'$ est vide ssi $F$ l'est. Prenons p.ex. \deleted{$\unsure{F|K}\unsure{\preceq}$} $\preceq$ \deleted{\unsure{égale}} \unsure{à|:} la relation d'ordre discrète, c'est OK. \deleted{$C7$ \unsure{a} est \unsure{tautologique}} $C6,C7$ tautologiques. \unsure{On voit donc|On veut dire} qu'il manque un axiome, pour assurer la validité du corollaire. On devine que ce sera la partie \unsure{a)|--)} du \unsure{Const|Cons} $4$ (p. 35), \unsure{jointe} à Const $7$. Donc pour l'instant je ne vais pas utiliser que le \unsure{préordre} $\ll$ soit un ordre.
% slanted
\margin{Aussi, cela résulte \unsure{aussitôt|aussi} \unsure{de la} \unsure{modification} de $C4$ \unsure{ancien} (page $47$).}

\deleted{\unsure{Je} songe à la propriété \unsure{suiv\textsuperscript{te}|suiv.}, qu'il faudrait avoir tôt ou tard :}

\deleted{\deleted{Soient $\unsure{Y|y},\unsure{X|x}\in M$ \unsure{avec} $Y\ll X$, et} supposons que \deleted{$Y'\leq Y$, $Y'\in\partial X$, $Y\cap\partial X=$} \deleted{$\unsure{\exists|\nexists}\,Y'\leq Y$, \unsure{alors}} $Y'\leq Y$ $(Y'\in M)$ tels que \deleted{$\unsure{\exists}\,Y'\ll\partial X$}}

\deleted{Soient $F$ figure, $F'$ sous-figure, $G\ll F$, $G'$ \inserted{$\leq G$} le \unsure{raffint|raffinement} de $F'$ induit par $G$ (cf plus bas). \deleted{supposons que} \inserted{Alors} $G'=\varnothing$. \deleted{Alors $G$ est} ssi $G$ disjoint de $F'$.}

Il faut développer la notion de \uline{\unsure{raffinement|raffinements} induit}. J'ai envie de \deleted{\unsure{démontrer}} \inserted{\unsure{prouver}} l'énoncé (quitte à rajouter des axiomes) :
